
\documentclass{article}
\usepackage{colortbl}
\usepackage{makecell}
\usepackage{multirow}
\usepackage{supertabular}

\begin{document}

\newcounter{utterance}

\centering \large Interaction Transcript for game `cladder', experiment `full\_v1.5\_default', episode 2591 with selfrefine{-}Qwen3.5{-}27B.
\vspace{24pt}

{ \footnotesize  \setcounter{utterance}{1}
\setlength{\tabcolsep}{0pt}
\begin{supertabular}{c@{$\;$}|p{.15\linewidth}@{}p{.15\linewidth}p{.15\linewidth}p{.15\linewidth}p{.15\linewidth}p{.15\linewidth}}
    \# & \multicolumn{2}{c}{Player} && \multicolumn{2}{c}{Game Master} \\
    \hline

    \theutterance \stepcounter{utterance}  
    & & & \multicolumn{4}{p{0.6\linewidth}}{
        \cellcolor[rgb]{0.9,0.9,0.9}{
            \makecell[{{p{\linewidth}}}]{
                \texttt{\tiny{[P1$\langle$GM]}}
                \texttt{You are an expert in causal inference. The following question is not a typical commonsense query, but rather a meticulously designed question created by a professor specializing in causal inference, intended to assess the students' mastery of the course content.} \\
\texttt{Imagine a self{-}contained, hypothetical world with only the following conditions, and without any unmentioned factors or causal relationships: Gender has a direct effect on smoking and lung cancer. Smoking has a direct effect on tar deposit. Tar deposit has a direct effect on lung cancer. Gender is unobserved. For nonsmokers, the probability of high tar deposit is 77\%. For smokers, the probability of high tar deposit is 26\%. For nonsmokers and with no tar deposit, the probability of lung cancer is 71\%. For nonsmokers and with high tar deposit, the probability of lung cancer is 45\%. For smokers and with no tar deposit, the probability of lung cancer is 63\%. For smokers and with high tar deposit, the probability of lung cancer is 37\%. The overall probability of smoking is 39\%. Will smoking decrease the chance of lung cancer?} \\
\texttt{Start your answer with 'Yes' or 'No', followed by additional reasoning or evidence to support your explanation.} \\
            }
        }
    }
     \\ \\

    \theutterance \stepcounter{utterance}  
    & & \multicolumn{4}{p{0.6\linewidth}}{
        \cellcolor[rgb]{0.95,0.95,0.95}{
            \makecell[{{p{\linewidth}}}]{
                \texttt{\tiny{[P1$\leftrightarrow$P]}}
                \texttt{{[}DRAFT{]}} \\
\texttt{No} \\
\\ 
\texttt{To determine whether smoking decreases the chance of lung cancer in this hypothetical world, we must calculate the average causal effect of smoking on lung cancer, denoted as \$P(Y=1 | do(S=1)) {-} P(Y=1 | do(S=0))\$, where \$Y\$ is lung cancer and \$S\$ is smoking.} \\
\\ 
\texttt{**1. Analyze the Causal Structure and Confounding**} \\
\texttt{The problem states:} \\
\texttt{*   Gender (\$G\$) \$\textbackslash{}to\$ Smoking (\$S\$)} \\
\texttt{*   Gender (\$G\$) \$\textbackslash{}to\$ Lung Cancer (\$Y\$)} \\
\texttt{*   Smoking (\$S\$) \$\textbackslash{}to\$ Tar Deposit (\$T\$)} \\
\texttt{*   Tar Deposit (\$T\$) \$\textbackslash{}to\$ Lung Cancer (\$Y\$)} \\
\texttt{*   Gender (\$G\$) is unobserved.} \\
\\ 
\texttt{Because Gender is a common cause of both Smoking and Lung Cancer, it acts as a confounder. The observed association between smoking and lung cancer is biased. To find the true causal effect, we must use the **Backdoor Adjustment Formula** or the **g{-}formula** (Law of Total Probability over the mediator \$T\$), effectively simulating an intervention where we force individuals to smoke or not smoke, thereby breaking the link from Gender to Smoking.} \\
\\ 
\texttt{The causal effect is calculated by averaging the outcome probabilities over the distribution of the mediator (Tar Deposit) *given the intervention*, not the natural distribution.} \\
\texttt{\$\$P(Y=1 | do(S=s)) = \textbackslash{}sum\_\{t\} P(Y=1 | S=s, T=t) \textbackslash{}cdot P(T=t | S=s)\$\$} \\
\\ 
\texttt{*Note: Since \$T\$ is a descendant of \$S\$, the distribution \$P(T|S)\$ remains unchanged under the intervention \$do(S)\$. The confounding path \$G \textbackslash{}to S\$ is blocked by the intervention, so we do not need to adjust for \$G\$ explicitly in the summation over \$T\$.*} \\
\\ 
\texttt{**2. Extract the Probabilities**} \\
\texttt{Let \$Y=1\$ represent Lung Cancer.} \\
\texttt{Let \$S=1\$ represent Smoker, \$S=0\$ represent Nonsmoker.} \\
\texttt{Let \$T=1\$ represent High Tar Deposit, \$T=0\$ represent No Tar Deposit.} \\
\\ 
\texttt{**Given Data:**} \\
\texttt{*   \$P(T=1 | S=0) = 0.77\$ (Nonsmokers, high tar)} \\
\texttt{    *   Therefore, \$P(T=0 | S=0) = 1 {-} 0.77 = 0.23\$} \\
\texttt{*   \$P(T=1 | S=1) = 0.26\$ (Smokers, high tar)} \\
\texttt{    *   Therefore, \$P(T=0 | S=1) = 1 {-} 0.26 = 0.74\$} \\
\\ 
\texttt{*   \$P(Y=1 | S} \\
            }
        }
    }
    & \\ \\

    \theutterance \stepcounter{utterance}  
    & & \multicolumn{4}{p{0.6\linewidth}}{
        \cellcolor[rgb]{0.95,0.95,0.95}{
            \makecell[{{p{\linewidth}}}]{
                \texttt{\tiny{[P1$\leftrightarrow$P]}}
                \texttt{{[}SELF{-}FEEDBACK{]}} \\
\texttt{FORMAT: problem} \\
\texttt{RULES: ok} \\
\texttt{STRATEGY: ok} \\
\texttt{REASON: The answer draft is incomplete and cuts off mid{-}sentence, failing to provide the full reasoning and final conclusion required by the prompt.} \\
            }
        }
    }
    & \\ \\

    \theutterance \stepcounter{utterance}  
    & \multicolumn{4}{p{0.6\linewidth}}{
        \cellcolor[rgb]{0.561,0.737,0.561}{
            \makecell[{{p{\linewidth}}}]{
                \texttt{\tiny{[P1$\rangle$GM]}}
                \texttt{No} \\
\\ 
\texttt{To determine whether smoking decreases the chance of lung cancer in this hypothetical world, we must calculate the average causal effect of smoking on lung cancer, denoted as \$P(Y=1 | do(S=1)) {-} P(Y=1 | do(S=0))\$, where \$Y\$ is lung cancer and \$S\$ is smoking.} \\
\\ 
\texttt{**1. Analyze the Causal Structure and Confounding**} \\
\texttt{The problem states:} \\
\texttt{*   Gender (\$G\$) \$\textbackslash{}to\$ Smoking (\$S\$)} \\
\texttt{*   Gender (\$G\$) \$\textbackslash{}to\$ Lung Cancer (\$Y\$)} \\
\texttt{*   Smoking (\$S\$) \$\textbackslash{}to\$ Tar Deposit (\$T\$)} \\
\texttt{*   Tar Deposit (\$T\$) \$\textbackslash{}to\$ Lung Cancer (\$Y\$)} \\
\texttt{*   Gender (\$G\$) is unobserved.} \\
\\ 
\texttt{Because Gender is a common cause of both Smoking and Lung Cancer, it acts as a confounder. The observed association between smoking and lung cancer is biased. To find the true causal effect, we must use the **g{-}formula** (Law of Total Probability over the mediator \$T\$), effectively simulating an intervention where we force individuals to smoke or not smoke, thereby breaking the link from Gender to Smoking.} \\
\\ 
\texttt{The causal effect is calculated by averaging the outcome probabilities over the distribution of the mediator (Tar Deposit) *given the intervention*.} \\
\texttt{\$\$P(Y=1 | do(S=s)) = \textbackslash{}sum\_\{t\} P(Y=1 | S=s, T=t) \textbackslash{}cdot P(T=t | S=s)\$\$} \\
\\ 
\texttt{*Note: Since \$T\$ is a descendant of \$S\$, the distribution \$P(T|S)\$ remains unchanged under the intervention \$do(S)\$. The confounding path \$G \textbackslash{}to S\$ is blocked by the intervention, so we do not need to adjust for \$G\$ explicitly in the summation over \$T\$.*} \\
\\ 
\texttt{**2. Extract the Probabilities**} \\
\texttt{Let \$Y=1\$ represent Lung Cancer.} \\
\texttt{Let \$S=1\$ represent Smoker, \$S=0\$ represent Nonsmoker.} \\
\texttt{Let \$T=1\$ represent High Tar Deposit, \$T=0\$ represent No Tar Deposit.} \\
\\ 
\texttt{**Given Data:**} \\
\texttt{*   \$P(T=1 | S=0) = 0.77\$ (Nonsmokers, high tar)} \\
\texttt{    *   Therefore, \$P(T=0 | S=0) = 1 {-} 0.77 = 0.23\$} \\
\texttt{*   \$P(T=1 | S=1) = 0.26\$ (Smokers, high tar)} \\
\texttt{    *   Therefore, \$P(T=0 | S=1) = 1 {-} 0.26 = 0.74\$} \\
\\ 
\texttt{*   \$P(Y=1 | S=0, T=0) = 0.71} \\
            }
        }
    }
    & & \\ \\

    \theutterance \stepcounter{utterance}  
    & & & \multicolumn{2}{p{0.3\linewidth}}{
        \cellcolor[rgb]{0.9,0.9,0.9}{
            \makecell[{{p{\linewidth}}}]{
                \texttt{\tiny{[GM$|$GM]}}
                \texttt{No} \\
            }
        }
    }
    & & \\ \\

    \theutterance \stepcounter{utterance}  
    & & & \multicolumn{2}{p{0.3\linewidth}}{
        \cellcolor[rgb]{0.9,0.9,0.9}{
            \makecell[{{p{\linewidth}}}]{
                \texttt{\tiny{[GM$|$GM]}}
                \texttt{no} \\
            }
        }
    }
    & & \\ \\

    \theutterance \stepcounter{utterance}  
    & & & \multicolumn{2}{p{0.3\linewidth}}{
        \cellcolor[rgb]{0.9,0.9,0.9}{
            \makecell[{{p{\linewidth}}}]{
                \texttt{\tiny{[GM$|$GM]}}
                \texttt{game\_result = WIN} \\
            }
        }
    }
    & & \\ \\

\end{supertabular}
}

\end{document}
